
\subsubsection{5.1 h-e1: Parameter Validity (PASS)}

\begin{table}[htbp]
\centering
\resizebox{\columnwidth}{!}{%
\begin{tabular}{llll}
\toprule
Metric & Threshold & Observed & Status \\
\midrule
NaN/Inf Rate & 0\% & 0.0\% & PASS \\
Magnitude Ratio & < $10\times$ & 0.$11\times$ & PASS \\
\bottomrule
\end{tabular}
}
\end{table}

Duality-derived parameters are numerically stable across all 100 test samples. The magnitude ratio of 0.$11\times$ indicates conservative parameter scales relative to attention outputs.

\subsubsection{5.2 h-m1: Reconstruction Error (FAIL)}

\begin{table}[htbp]
\centering
\resizebox{\columnwidth}{!}{%
\begin{tabular}{ll}
\toprule
Metric & Value \\
\midrule
Mean Duality Error & 91.45 \\
Mean Random Error & 89.54 \\
Error Reduction & $-2.13$\% (worse) \\
P-value & < 0.0001 \\
Cohen's d & $-4.56$ (large negative) \\
\bottomrule
\end{tabular}
}
\end{table}

Duality initialization produces 2.13\% higher reconstruction error than random initialization. The effect is statistically significant (p < 0.0001) with a large negative effect size (Cohen's d = $-4.56)$.

\subsubsection{5.3 Per-Layer Analysis}

\begin{table*}[htbp]
\centering
\resizebox{\dimexpr 2\columnwidth+\columnsep\relax}{!}{%
\begin{tabular}{lllll}
\toprule
Layer & Duality Error & Random Error & Reduction & Cohen's d \\
\midrule
0 & 79.75 & 77.74 & $-2.58$\% & $-5.02$ \\
1 & 93.98 & 92.16 & $-1.98$\% & $-5.50$ \\
2 & 108.35 & 106.70 & $-1.55$\% & $-5.19$ \\
3 & 104.67 & 103.04 & $-1.57$\% & $-5.97$ \\
4 & 98.18 & 96.35 & $-1.89$\% & $-4.15$ \\
5 & 96.58 & 94.85 & $-1.82$\% & $-5.92$ \\
6 & 96.29 & 94.53 & $-1.86$\% & $-5.12$ \\
7 & 91.16 & 89.18 & $-2.22$\% & $-4.71$ \\
8 & 84.44 & 82.31 & $-2.58$\% & $-5.85$ \\
9 & 87.10 & 85.06 & $-2.40$\% & $-5.38$ \\
10 & 79.79 & 77.64 & $-2.77$\% & $-5.31$ \\
11 & 77.11 & 74.91 & $-2.94$\% & $-4.84$ \\
\bottomrule
\end{tabular}
}
\end{table*}

Duality initialization is consistently worse across all 12 BERT layers, with Cohen's d ranging from $-4.15$ to $-5.97$ (all large negative effects). Layer 11 shows the largest degradation ($-2.94$\%), while Layer 2 shows the smallest ($-1.55$\%).

\subsubsection{5.4 Causal Chain Verification}

\begin{table}[htbp]
\centering
\resizebox{\columnwidth}{!}{%
\begin{tabular}{lll}
\toprule
Step & Description & Status \\
\midrule
1 & Duality equations $\rightarrow$ valid SSM parameters & VERIFIED \\
2 & Duality parameters $\rightarrow$ lower reconstruction error & FALSIFIED \\
3 & Lower error $\rightarrow$ better task performance & BLOCKED \\
\bottomrule
\end{tabular}
}
\end{table}

The causal chain breaks at Step 2. While duality equations produce valid parameters, these parameters do not capture attention structure for output-level reconstruction.
